\section{Introduction}

Modern data-serving systems often place many storage-engine tenants on one
host. Each tenant may own a separate \rocksdb instance, data directory,
configuration, and foreground workload, yet their flushes and compactions draw
from the same finite background CPU and storage bandwidth. Per-instance controls are
useful, but they do not decide how a fixed host budget should be divided
across competing instances.

This placement decision becomes a performance problem when demand drifts. Equal
static allocation leaves capacity assigned to quiet tenants; frozen skew becomes
stale when the high-pressure set moves; and RocksDB's per-instance auto-tuned
rate limiter adapts within a configured ceiling but cannot move budget across
instances. The resulting compaction-debt mismatch appears as high write-tail
latency and lower completion rates for tenants currently creating debt.

We study the host-level placement decision directly. Under a fixed aggregate
background compaction/I/O budget, can an online controller use runtime
observations to reallocate per-tenant budgets across co-located \rocksdb
tenants, improving the tenants currently under write pressure without oracle
labels, workload forecasts, or extra host capacity? The question is
not whether more compaction capacity helps, but whether the same capacity can
be placed better than fixed fair sharing.

Several close directions leave this placement point open. General
multi-tenant storage controllers enforce lower-layer I/O shares or SLOs, while
recent LSM-oriented systems either change a single engine, move flush or
compaction into a disaggregated service or DPU path, or build isolation into a
different multi-tenant engine boundary~\cite{gulati2010mclock,
thereska2013ioflow,yu2024caaslsm,ding2025dflush,wang2026flexengine,
griggs2026deltafairsharing}. SAKI instead targets a common operational shape:
many separate, unmodified \rocksdb instances on one host, with only a conserved
background budget to place across them as demand drifts.

SAKI is an external controller for this placement problem. It observes each
tenant over a control window, ranks tenants using offered write demand and
compaction-pressure symptoms, and applies the next-window budget assignment
while conserving the aggregate budget. This conserved rank-and-assign framework
separates the placement invariant from the ranking key. Its key condition is
demand observability: when offered demand separates currently high-pressure
tenants from the rest, demand-aware ranking keys track bounded drift within
one control window, while pressure-only ranking can chase stale debt. SAKI sits
outside RocksDB; in the continuous harness, it changes runtime rate-limiter
budgets through \ratecall.

The evaluation follows an attribution chain. It first asks whether SAKI improves
the target tenants under the same aggregate budget, using both epoch-level and
continuous \rocksdb experiments. It then tests alternative explanations,
including static skew, tenant-local autotuning, stale-debt scoring, lower-layer
I/O weights, and foreground throttling. Finally, ranking-margin replay,
ranking-key perturbations, workload-matrix experiments, collateral measurements,
liveness checks, and public-trace audits define where the evidence applies.

In a 16-tenant epoch-level workload, five trials show that SAKI reduces
high-pressure write P99 by 11.7\%, raises high-pressure throughput by 38.1\%,
and lowers compaction write bytes by 21.0\% relative to equal static allocation,
with roughly neutral total throughput. In a continuous embedded \rocksdb harness
with runtime actuation, five trials show that SAKI reduces high-pressure write
P99 by 27.0\%, reduces P999 by 47.1\%, raises high-pressure throughput by
26.1\%, and lowers normalized bytes/write by 11.9\%. The tradeoff is visible:
lower-pressure tenants may absorb part of the cost, so SAKI is a targeted
reallocation policy.

The paper makes five contributions:
\begin{enumerate}[leftmargin=1.2em,itemsep=0pt,topsep=2pt]
\item It frames multi-tenant \rocksdb compaction as a fixed-budget host-level
reallocation problem, distinct from single-instance compaction-policy design.
\item It presents SAKI, a conserved rank-and-assign controller that sets
per-tenant runtime budgets from previous-window observations without oracle
labels, workload forecasts, or tenant application changes.
\item It identifies the demand-observability condition that lets a
demand-aware ranking key track bounded drift, and contrasts it with the
stale-debt lag of pressure-only ranking.
\item It provides epoch-level and continuous evidence that online reallocation
improves high-pressure write latency and throughput under bounded drift while
keeping total throughput close to static allocation.
\item It separates SAKI's effect from fixed skew, per-instance local autotuning,
ranking-key choice, workload variation, and lower-layer I/O control, while
making collateral costs and validation boundaries explicit.
\end{enumerate}

Together, these results show that shared host-level compaction budgeting is a
practical control point for multi-tenant \rocksdb write performance under
bounded drift.
